\section{Nonvanishing with fewer nodes}\label{sec:rank}
The last algebraic input is full row rank of the matrix $V$. A dimension count
requires $nJ\geq M$, but it does not prove that the rows are independent.
The criterion below comes close to that count and is sufficient for all the
finite comparisons between $499$ and $875$. The proof uses a divided-Frobenius basis to retain information lost by direct reduction modulo a prime.

Direct reduction modulo an inert prime loses information: the polynomial
$z^p-w$ vanishes at every reduced node. The remedy is to divide out the
powers of $p$ whose presence we already know. After this division, a hypothetical
row relation gives a polynomial in two variables with many prescribed local
zeros. A differential operator converts those conditions into too many zeros
of a nonzero polynomial in one variable. The two main points are to prove that
this last polynomial is nonzero and to compare its degree with its orders of
vanishing.

\begin{theorem}[Sufficient rank criterion]\label{thm:simple-rank}
Let $p$ be an inert prime, and suppose that the nodes $z_1,\ldots,z_n$ have
distinct residues modulo $p\OO_D$. Let $R=KS$ and
\[
 K<p<R\leq p^2,\qquad 1\leq J\leq R.
\]
If $n\geq2S$ and
\begin{equation}\label{eq:simple-rank-ineq}
 n\left(J-\frac{RJ}{p^2}-\frac Rp\right)>
 \frac{(R+p)^2}{2K},
\end{equation}
then the matrix $V$ in \eqref{eq:jet-matrix} has full row rank.
\end{theorem}

The leading terms of \eqref{eq:simple-rank-ineq} are $nJ>R^2/(2K)$, close
to the necessary inequality $nJ\geq M=R(S+1)/2$. The errors are small when
$p/R$ and $R/p^2$ are small. For our application the choice of $p$ must also
allow distinct node residues and guarantee that at least one candidate prime
is inert. These requirements will be checked together in
Section~\ref{sec:completion}.

\subsection{Dividing out the Frobenius factors}
Put $g=z^p-w$ and $d=\lfloor(R-1)/p\rfloor<p$. The integral basis change
\eqref{eq:frobenius-basis} gives the rows
\begin{equation}\label{eq:g-basis}
 g^m z^a w^s,\qquad 0\leq m\leq d,\quad 0\leq a<p,\quad a+Ks<R-mp.
\end{equation}
At node $i$ write
\[
 g(z_i+t,\bar z_i)=pU_i(t)+t^p,
 \quad
 U_i(t)=\frac{z_i^p-\bar z_i}{p}
       +\sum_{v=1}^{p-1}\frac1p\binom pv z_i^{p-v}t^v\in\OO_D[t].
\]
Direct reduction would replace $(pU_i(t)+t^p)^m$ by $t^{mp}$, discarding
$U_i$ and making layer $m$ vanish in every column of order less than $mp$.
Instead, divide each row in layer $m$ by $p^m$ and multiply each column of order
$j=qp+v$, $0\leq v<p$, by $p^q$. Call the resulting matrix $\widetilde V$.
Its entries are integral. Indeed, in the binomial expansion of
$(pU_i+t^p)^m$, the term indexed by $q'$ can contribute to order $j$ only
when $q'\leq q$; after rescaling its coefficient contains $p^{q-q'}$.
Modulo $p$, only the term $q'=q$ survives.

Write $k=\OO_D/p\OO_D$, $\xi_i=z_i\bmod p$, and
$u_i(z)=(U_i\bmod p)(z-\xi_i)$. Thus, when $q\leq m$, the reduced entry is
\begin{equation}\label{eq:scaled-entry}
 \binom mq[t^v]\bigl(u_i(\xi_i+t)^{m-q}(\xi_i+t)^a\xi_i^{ps}\bigr).
\end{equation}
When $q>m$, all terms are divisible by $p$, and the reduced entry is zero.
The basis change and all scalings are invertible in characteristic zero.
It therefore suffices to prove that $\widetilde V\bmod p$ has full row rank
over $k$. These scalings are used only for nonvanishing; the determinant
comparison continues to concern the original matrix $V$.

Suppose, to the contrary, that the reduced rows satisfy a nontrivial relation
with coefficients $c_{m,a,s}\in k$. Package that relation as
\begingroup\postdisplaypenalty=10000
\begin{equation}\label{eq:rank-F}
 F(z,u)=\sum_{m=0}^d H_m(z)u^m,\qquad
 H_m(z)=\sum_{a,s}c_{m,a,s}z^{a+ps}.
\end{equation}
\endgroup
Within each layer the exponents $a+ps$ are distinct, since $0\leq a<p$.
Consequently $F\ne0$. Write $e=\deg_uF\leq d<p$. The triangular support gives
\begin{equation}\label{eq:rank-degree}
 \deg H_m<E_m:=\frac{p(R-mp)}K,
\end{equation}
because
$a+ps=(p/K)(a+Ks)-(p/K-1)a<p(R-mp)/K$.
The importance of this bound is its constant drop:
$E_m-E_{m+1}=p^2/K>p$.

To express the local conditions on $F$, put
\[
 L_l=\min\bigl(p,\max(0,J-lp)\bigr).
\]
There are exactly $L_l$ derivative orders in the block $j=lp+v$.
In characteristic $p$, $(\xi_i+t)^{ps}\equiv\xi_i^{ps}\pmod{t^p}$.
We may therefore replace $\xi_i^{ps}$ by $(\xi_i+t)^{ps}$ when taking any
of the coefficients $[t^v]$ with $v<L_l$. Summing \eqref{eq:scaled-entry}
against the relation coefficients gives
\begin{equation}\label{eq:local-hasse-vanish}
 \ord_{z=\xi_i}\bigl((\partial_u^{[l]}F)(z,u_i(z))\bigr)\geq L_l
 \qquad(i=1,\ldots,n;\ l\geq0).
\end{equation}
Here $\partial_u^{[l]}$ denotes the Hasse derivative in $u$. Thus a relation
among the scaled rows has become a family of vanishing conditions along
polynomial curves $u=u_i(z)$.

\subsection{One differential operator for all nodes}
Differentiating $pU_i(t)=(z_i+t)^p-\bar z_i-t^p$ before reducing modulo $p$
gives
\begin{equation}\label{eq:ui-prime}
 u_i'(z)=z^{p-1}-(z-\xi_i)^{p-1}.
\end{equation}
The first term is independent of the node, and the second vanishes to order
$p-1$ there. This suggests using the common operator
\[
 \mathcal D=\partial_z+z^{p-1}\partial_u
\]
in place of differentiation along each separate curve. Its key property is
\begin{equation}\label{eq:Dp-identity}
 \mathcal D^p=-\partial_u\qquad\hbox{on }k(z)[u].
\end{equation}
Indeed, the $p$th iterate of a derivation in characteristic $p$ is a derivation.
It sends $z$ to zero and $u$ to $(p-1)!=-1$, which determines it on $k(z)[u]$.

\begin{lemma}\label{lem:D-cyclic}
The polynomials $F,\mathcal DF,\ldots,\mathcal D^eF$ are linearly independent
over $k(z)$. In particular
\begin{equation}\label{eq:WF-def}
 W_F(z)=\det\bigl([u^m]\mathcal D^jF\bigr)_{0\leq j,m\leq e}
\end{equation}
is nonzero, and
\begin{equation}\label{eq:WF-degree}
 \deg W_F\leq\sum_{m=0}^e E_m-\frac{e(e+1)}2.
\end{equation}
\end{lemma}
\begin{proof}
If these polynomials were dependent, let $r\leq e$ be the first index of a
dependence. Their span up to index $r-1$ would be stable under $\mathcal D$:
applying $\mathcal D$ to a rational coefficient differentiates that coefficient,
and the image of the last spanning vector is in the same span by assumption.
By \eqref{eq:Dp-identity} this space is also stable under $\partial_u$.
It would then contain
$F,\partial_uF,\ldots,\partial_u^eF$. Since $e<p$, these have distinct
$u$-degrees $e,e-1,\ldots,0$ and nonzero leading coefficients. They are
independent, contradicting the dimension $r\leq e$.

For the degree bound, induction using
\[
 [u^m]\mathcal DG=([u^m]G)' +(m+1)z^{p-1}[u^{m+1}]G
\]
shows that $\deg([u^m]\mathcal D^jF)\leq E_m-j$. The first term lowers
degree by one; the second has degree at most
$E_{m+1}-j+p-1<E_m-j-1$. In every determinant term each column index $m$
and row index $j$ occurs once. Adding these degree bounds proves
\eqref{eq:WF-degree}.
\end{proof}

The use of $e<p$ here is essential: it prevents the successive $u$-derivatives
from vanishing prematurely in characteristic $p$. The identity
$\mathcal D^p=-\partial_u$ is what connects this elementary fact to
differentiation in $z$.

\subsection{Counting zeros of the coefficient determinant}
Fix a node $i$. Translation from the basis $1,u,\ldots,u^e$ to
$1,u-u_i(z),\ldots,(u-u_i(z))^e$ is unitriangular. It leaves $W_F$ unchanged
and replaces its entries by
\[
 B^{(i)}_{j,l}(z)=(\partial_u^{[l]}\mathcal D^jF)(z,u_i(z)).
\]
We claim that
\begin{equation}\label{eq:Bi-order-bound}
 \ord_{\xi_i}B^{(i)}_{j,l}\geq\max(0,L_l-j).
\end{equation}
To see this, let $\mathcal D_i=\partial_z+u_i'(z)\partial_u$, which differentiates
along the curve. Equation~\eqref{eq:ui-prime} gives
$\mathcal D-\mathcal D_i=(z-\xi_i)^{p-1}\partial_u$. An induction, in which
each differentiation lowers the order of the error by at most one, gives
\[
 \mathcal D^jG\equiv\mathcal D_i^jG
       \pmod{(z-\xi_i)^{p-j}k[z,u]}\qquad(0\leq j<p).
\]
The Hasse derivatives in $u$ commute with both operators: their coefficients
depend only on $z$, and ordinary and Hasse $u$-derivatives commute.
After specializing $u=u_i(z)$, the corresponding entry is therefore congruent
to the $j$th ordinary derivative of
$(\partial_u^{[l]}F)(z,u_i(z))$, modulo $(z-\xi_i)^{p-j}$.
Its order is at least $\max(0,L_l-j)$ by \eqref{eq:local-hasse-vanish}.
Since $L_l\leq p$, the error has at least that order as well. This proves the claim.

We can now finish the proof by factoring powers of $z-\xi_i$ out of rows
or columns. Put $Q=\lfloor J/p\rfloor$. If $e<Q$, then $L_l=p$ for all
$0\leq l\leq e$. Every entry in row $j$ of $B^{(i)}$ has order at least
$p-j$, so
\[
 \ord_{\xi_i}W_F\geq\sum_{j=0}^e(p-j)
 =(e+1)(p-e/2)>(e+1)p/2.
\]
The residues are distinct and $n\geq2S=2R/K$. Summing this bound over the nodes
therefore gives more than $(e+1)pR/K$ zeros counted with multiplicity,
whereas \eqref{eq:WF-degree} bounds the degree by at most that number.
This contradicts $W_F\ne0$.

If $e\geq Q$, every entry in column $l$ has order at least
$\max(0,L_l-e)$. Factoring column by column and using
$\sum_{l=0}^Q L_l=J$ gives
\begin{equation}\label{eq:rank-local-total}
 \ord_{\xi_i}W_F\geq\sum_{l=0}^Q(L_l-e)
 =J-(Q+1)e\geq J-\frac{RJ}{p^2}-\frac Rp.
\end{equation}
For the degree, write $R=dp+s_0$, where $1\leq s_0\leq p$. Then
\[
 2p\sum_{m=0}^d(R-mp)
 =(R-s_0+p)(R+s_0)
 =R^2+pR+s_0(p-s_0)<(R+p)^2.
\]
All $E_m$ are positive. Hence
\[
 \deg W_F\leq\sum_{m=0}^e E_m\leq\sum_{m=0}^d E_m
 <\frac{(R+p)^2}{2K}.
\]
Summing \eqref{eq:rank-local-total} over the nodes and applying
\eqref{eq:simple-rank-ineq} again makes the total order of vanishing exceed
the degree. Both cases are impossible. The reduced matrix has full row rank,
and so does $V$, proving Theorem~\ref{thm:simple-rank}.
